\documentclass[12pt]{article}
\usepackage{amsmath,amssymb}

\begin{document}
\section*{{2020 AIME I Problems/Problem 12}}
Source: AoPS Wiki

\subsection*{Solution}
As usual, denote $v_p(n)$ the highest power of prime $p$ that divides $n$ . Lifting the Exponent shows that \[3=v_3(149^n-2^n) = v_3(n) + v_3(147) = v_3(n)+1\] so thus, $3^2$ divides $n$ . It also shows that \[7=v_7(149^n-2^n) = v_7(n) + v_7(147) = v_7(n)+2\] so thus, $7^5$ divides $n$ . Now, setting $n = 4c$ (necessitated by $149^n \equiv 2^n \pmod 5$ in order to set up LTE), we see \[v_5(149^{4c}-2^{4c}) = v_5(149^{4c}-16^{c})\] and since $149^{4} \equiv 1 \pmod{25}$ and $16^1 \equiv 16 \pmod{25}$ then $v_5(149^{4c}-2^{4c})=v_5(149^4-16)+v_5(c)=1+v_5(c)$ meaning that we have that by LTE, $5^4 | c$ and $4 \cdot 5^4$ divides $n$ . Since $3^2$ , $7^5$ and $4\cdot 5^4$ all divide $n$ , the smallest value of $n$ working is their LCM, also $3^2 \cdot 7^5 \cdot 4 \cdot 5^4 = 2^2 \cdot 3^2 \cdot 5^4 \cdot 7^5$ . Thus the number of divisors is $(2+1)(2+1)(4+1)(5+1) = \boxed{270}$ . ~kevinmathz clarified by another user notation note from another user

\end{document}
